\heading{经济学原理（II）-模拟题1-期中}
\noteinfo{题目和答案由 AI 依据样题与第 1--6 讲内容生成；本卷较样题更强调概念辨析、口径判断与多步推导。}

\instructions{建议答题时间 120 分钟。若无特别说明，百分数按小数形式计算后再换算；计算题请写出关键中间步骤。}

\textbf{一、单项选择题}

\begin{question}
在一条向外凸出的生产可能性边界上，如果经济从左上方移动到右下方，意味着
\begin{choiceoptions}[2]
  \optionitem{生产大炮的机会成本递减}
  \optionitem{生产黄油的机会成本恒定}
  \optionitem{生产大炮的边际机会成本递增}
  \optionitem{资源配置越来越有效率}
\end{choiceoptions}
\begin{solution}
向外凸出的生产可能性边界表示随着某种产品产量继续增加，需要放弃的另一种产品越来越多，因此边际机会成本递增。
\anstext{C}
\end{solution}
\end{question}

\begin{question}
下列哪一项在 2026 年\textbf{不计入}中国 GDP？
\begin{choiceoptions}[2]
  \optionitem{中国家庭购买当年国内新建商品房}
  \optionitem{中国企业购买当年生产的国产机床}
  \optionitem{中国居民购买一辆 2024 年生产的二手汽车}
  \optionitem{房产中介在二手房交易中收取的佣金}
\end{choiceoptions}
\begin{solution}
新房属于投资，机床属于投资，中介佣金属于当期服务；二手车本身不是本期生产，不计入当期 GDP。
\anstext{C}
\end{solution}
\end{question}

\begin{question}
关于平均税率与边际税率，下列说法正确的是
\begin{choiceoptions}[2]
  \optionitem{平均税率更能反映劳动供给激励}
  \optionitem{边际税率更能反映额外挣一单位收入所面临的激励变化}
  \optionitem{总额税的边际税率一定最高}
  \optionitem{累进税制下一定有边际税率低于平均税率}
\end{choiceoptions}
\begin{solution}
行为决策往往取决于“再多做一点”的收益与成本，因此边际税率更能反映激励。总额税的边际税率通常为零；累进税制下边际税率常高于平均税率。
\anstext{B}
\end{solution}
\end{question}

\begin{question}
若一国境内外资企业较多，且本国居民从国外获得的要素收入较少，则通常有
\begin{choiceoptions}[2]
  \optionitem{GNP $>$ GDP}
  \optionitem{GDP $>$ GNP}
  \optionitem{GDP 一定等于 GNP}
  \optionitem{GDP 与 GNP 无法比较}
\end{choiceoptions}
\begin{solution}
GDP 是地域口径，GNP 是居民口径。若境内有较多外国资本获取收益而本国居民来自国外的收入较少，则净要素收入为负，故 GNP 小于 GDP。
\anstext{B}
\end{solution}
\end{question}

\begin{question}
下列哪一项最可能使 CPI 高估真实生活费用上涨？
\begin{choiceoptions}[2]
  \optionitem{居民在相对价格变化后用更便宜的替代品替换原商品}
  \optionitem{政府购买增加}
  \optionitem{出口增加}
  \optionitem{人口增长放缓}
\end{choiceoptions}
\begin{solution}
固定消费篮子没有及时反映替代行为，这就是 CPI 的替代偏差。
\anstext{A}
\end{solution}
\end{question}

\begin{question}
在索罗模型中，若实际人均投资小于维持资本密度不变所需的人均投资，则
\begin{choiceoptions}[2]
  \optionitem{人均资本上升}
  \optionitem{人均资本下降}
  \optionitem{人均资本不变但总资本上升}
  \optionitem{总人口下降}
\end{choiceoptions}
\begin{solution}
稳态条件是 $sf(k)=(n+\delta)k$。若左边更小，则人均资本会下降。
\anstext{B}
\end{solution}
\end{question}

\begin{question}
若政府从预算平衡转向预算赤字，其他条件不变，则在可贷资金市场中通常会出现
\begin{choiceoptions}[2]
  \optionitem{储蓄供给右移，真实利率下降，投资上升}
  \optionitem{储蓄供给左移，真实利率上升，投资下降}
  \optionitem{投资需求左移，真实利率下降，投资下降}
  \optionitem{投资需求右移，真实利率上升，投资上升}
\end{choiceoptions}
\begin{solution}
预算赤字降低公共储蓄，从而降低国民储蓄，表现为可贷资金供给左移，真实利率上升，投资被部分挤出。
\anstext{B}
\end{solution}
\end{question}

\begin{question}
若预期通胀上升而真实利率不变，则依据费雪效应
\begin{choiceoptions}[2]
  \optionitem{名义利率下降}
  \optionitem{名义利率上升}
  \optionitem{真实利率下降到零}
  \optionitem{储蓄与投资必然同时上升}
\end{choiceoptions}
\begin{solution}
费雪方程近似为 $i\approx r+\pi^e$。当预期通胀上升、真实利率不变时，名义利率上升。
\anstext{B}
\end{solution}
\end{question}

\begin{question}
分散化能够显著降低的主要是
\begin{choiceoptions}[2]
  \optionitem{系统性风险}
  \optionitem{通货膨胀风险}
  \optionitem{非系统性风险}
  \optionitem{利率风险中的共同宏观部分}
\end{choiceoptions}
\begin{solution}
分散化主要消除个别资产特有的波动，即非系统性风险；系统性风险无法仅靠持有更多资产消除。
\anstext{C}
\end{solution}
\end{question}

\begin{question}
半强式有效市场假说最恰当的表述是
\begin{choiceoptions}[2]
  \optionitem{任何信息都不会影响价格}
  \optionitem{公开信息会迅速反映在资产价格中}
  \optionitem{市场价格永不偏离基本面}
  \optionitem{只有内幕信息是无价值的}
\end{choiceoptions}
\begin{solution}
半强式有效市场假说强调公开信息被迅速反映在价格中，但并不等于价格绝不偏离基本面。
\anstext{B}
\end{solution}
\end{question}

\textbf{二、简答题}

\begin{question}
为什么转移支付不计入政府购买 $G$，但它又可能通过其他渠道影响 GDP？

\begin{solution}
转移支付不计入 $G$，因为它不是对本期最终产品和服务的购买，而只是收入在家庭、政府等主体之间的再分配，比如养老金、失业补助等。GDP 统计的是当期生产而不是收入转手本身，所以转移支付不直接进入支出法中的 $G$。但转移支付会改变居民可支配收入，进而影响消费 $C$；若消费因此上升，就会通过支出法间接影响 GDP。也就是说，\textbf{转移支付不直接计入 GDP，但会通过影响行为间接改变 GDP。}
\end{solution}
\end{question}

\begin{question}
比较 GDP 平减指数与 CPI 的覆盖范围，并说明为什么两者的变动幅度可能不同。

\begin{solution}
GDP 平减指数覆盖的是\textbf{国内生产}的最终产品和服务，因此政府购买、投资品价格都会影响它；CPI 覆盖的是\textbf{居民消费}的固定篮子，因此进口消费品也会影响它。两者变动幅度可能不同，原因在于：第一，覆盖对象不同；第二，CPI 使用固定篮子，容易出现替代偏差；第三，GDP 平减指数随当期产出结构变化而变化。例如进口消费品价格大涨时，CPI 可能显著上升，而 GDP 平减指数未必同步上升。
\end{solution}
\end{question}

\begin{question}
在索罗模型中，为什么提高储蓄率可以提高长期产出水平，却不一定提高长期增长率？

\begin{solution}
提高储蓄率会增加投资，从而提高资本积累速度，使经济向更高的人均资本和更高的人均产出稳态过渡。因此长期\textbf{水平}会提高。但由于资本存在边际报酬递减，单靠提高储蓄率不能永久维持更高的人均增长率；当经济到达新的稳态后，人均增长再次回到零（若无外生技术进步）。所以更高的储蓄率带来的是\textbf{向更高稳态的过渡增长}，而不是永久提高长期增长率。
\end{solution}
\end{question}

\begin{question}
解释“分散化降低非系统性风险”这一命题，为什么它在现实中不是“风险可以完全消失”的同义词？

\begin{solution}
非系统性风险来自个别企业或个别资产特有事件，例如管理失误、产品失败、单个行业冲击等。若将多种相关性不完全的资产放入一个组合，个别波动可以部分相互抵消，因此总风险下降。可是系统性风险来自宏观因素，如衰退、利率冲击、通胀或金融危机，这类风险会同时影响大量资产，无法仅靠“多持有几只资产”消除。因此分散化降低的是\textbf{可分散风险}，不是让总风险必然变成零。
\end{solution}
\end{question}

\textbf{三、综合计算题}

\begin{question}
设某经济体在 2026 年发生如下交易（单位：亿元）：

\begin{itemize}[leftmargin=2em,itemsep=0.2em,topsep=0.2em]
  \item 居民购买国产食品 100、理发服务 50、进口笔记本电脑 40；
  \item 一套当年新建住房以 200 的价格卖给家庭；
  \item 一套旧房以 300 的价格转手，中介收取佣金 10；
  \item 企业购买国产机器 60，当年新增存货 30；
  \item 政府购买当年服务 90，并向居民发放养老金 40；
  \item 本国出口 70。
\end{itemize}

\begin{enumerate}[label=(\arabic*),leftmargin=2.5em]
  \item 计算消费 $C$、投资 $I$、政府购买 $G$、净出口 $NX$；
  \item 计算该年 GDP；
  \item 说明为什么养老金和旧房本身的交易口径处理不同。
\end{enumerate}

\begin{solution}
按支出法口径：
\[
C=100+50+40+10=200.
\]
其中中介佣金 10 属于本期服务，计入消费；新房不计入消费。

投资包括新建住房、企业设备投资和存货投资，因此
\[
I=200+60+30=290.
\]

政府购买只包括政府对最终产品和服务的购买，
\[
G=90.
\]
养老金属于转移支付，不计入 $G$。

净出口为
\[
NX=X-M=70-40=30.
\]

故 GDP 为
\[
Y=C+I+G+NX=200+290+90+30=610.
\]
\ansmath{C=200,\quad I=290,\quad G=90,\quad NX=30,\quad GDP=610.}

养老金不计入 GDP，是因为它不是对本期生产的支付；旧房本身也不计入 GDP，是因为它不是本期生产。但旧房交易中的中介佣金是本期新增服务，所以要计入 GDP。两者都不属于“当期生产的最终产品本身”，但\textbf{中介服务}属于本期生产。
\end{solution}
\end{question}

\begin{question}
某国五等份家庭收入份额依次为 $4\%$、$8\%$、$13\%$、$22\%$、$53\%$。

\begin{enumerate}[label=(\arabic*),leftmargin=2.5em]
  \item 写出洛伦兹曲线的 6 个关键点（含原点与终点）；
  \item 用分段线性近似计算基尼系数；
  \item 若政府通过税收和转移支付使收入份额变为 $8\%$、$12\%$、$13\%$、$22\%$、$45\%$，重新计算基尼系数并判断不平等如何变化。
\end{enumerate}

\begin{solution}
原分配下的累计人口占比与累计收入占比分别为
\[
(0,0),\ (0.2,0.04),\ (0.4,0.12),\ (0.6,0.25),\ (0.8,0.47),\ (1,1).
\]

洛伦兹曲线下方面积用梯形面积近似：
\[
A_L=\sum_{i=1}^{5}\frac{0.2\,(y_{i-1}+y_i)}{2}
=0.004+0.016+0.037+0.072+0.147=0.276.
\]
故基尼系数为
\[
G=1-2A_L=1-2\times 0.276=0.448.
\]
\ansmath{G_{\text{原}}=0.448.}

调整后的累计收入占比分别为
\[
(0,0),\ (0.2,0.08),\ (0.4,0.20),\ (0.6,0.33),\ (0.8,0.55),\ (1,1).
\]
此时面积为
\[
A'_L=0.008+0.028+0.053+0.088+0.166=0.344.
\]
故新的基尼系数
\[
G'=1-2A'_L=1-2\times 0.344=0.336.
\]
\ansmath{G_{\text{新}}=0.336.}

因为基尼系数从 $0.448$ 降到 $0.336$，说明收入分配更趋平等，不平等程度下降。
\end{solution}
\end{question}

\begin{question}
某经济体只生产三种商品 $A,B,C$，数据如下（数量，价格）：

\begin{center}
\begin{tabular}{c|ccc}
年份 & $A$ & $B$ & $C$\\ \hline
2025 & $(10,2)$ & $(8,5)$ & $(6,4)$\\
2026 & $(12,3)$ & $(7,6)$ & $(8,5)$
\end{tabular}
\end{center}

另知代表性消费者的固定消费篮子为：$2$ 单位 $A$，$1$ 单位 $B$，$3$ 单位 $C$。

\begin{enumerate}[label=(\arabic*),leftmargin=2.5em]
  \item 计算 2025 与 2026 年名义 GDP；
  \item 以 2025 年为基期，计算 2026 年真实 GDP 与真实增长率；
  \item 计算 2026 年 GDP 平减指数；
  \item 以 2025 年为基期，计算 2026 年 CPI。
\end{enumerate}

\begin{solution}
2025 年名义 GDP：
\[
NGDP_{2025}=10\times 2+8\times 5+6\times 4=84.
\]
2026 年名义 GDP：
\[
NGDP_{2026}=12\times 3+7\times 6+8\times 5=118.
\]

以 2025 年为基期，2026 年真实 GDP 为
\[
RGDP_{2026}^{(2025)}=12\times 2+7\times 5+8\times 4=91.
\]
由于 2025 年基期下真实 GDP 就是 84，所以真实增长率为
\[
\frac{91}{84}-1=\frac{7}{84}\approx 8.33\%.
\]

2026 年 GDP 平减指数为
\[
\text{Deflator}_{2026}=\frac{NGDP_{2026}}{RGDP_{2026}^{(2025)}}\times 100
=\frac{118}{91}\times 100\approx 129.67.
\]

固定消费篮子在 2025 年成本为
\[
2\times 2+1\times 5+3\times 4=21.
\]
在 2026 年成本为
\[
2\times 3+1\times 6+3\times 5=27.
\]
故 2026 年 CPI 为
\[
CPI_{2026}=\frac{27}{21}\times 100\approx 128.57.
\]

\ansmath{NGDP_{2025}=84,\quad NGDP_{2026}=118,}
\ansmath{RGDP_{2026}^{(2025)}=91,\quad g_{\text{real}}\approx 8.33\%,}
\ansmath{\text{Deflator}_{2026}\approx 129.67,\quad CPI_{2026}\approx 128.57.}
\end{solution}
\end{question}

\begin{question}
设索罗模型中生产函数为
\[
y=k^{1/3},
\]
储蓄率 $s=0.16$，人口增长率 $n=0.01$，折旧率 $\delta=0.03$。假定初始人均资本为 $k_0=27$。

\begin{enumerate}[label=(\arabic*),leftmargin=2.5em]
  \item 计算当前人均产出、实际人均投资与维持资本密度不变所需投资；
  \item 判断人均资本是上升还是下降；
  \item 若采用离散时间更新式
  \[
  k_{t+1}=\frac{(1-\delta)k_t+s f(k_t)}{1+n},
  \]
  计算 $k_1$；
  \item 求稳态人均资本 $k^*$；
  \item 求黄金规则储蓄率。
\end{enumerate}

\begin{solution}
因为
\[
y_0=f(k_0)=27^{1/3}=3,
\]
所以实际人均投资为
\[
s f(k_0)=0.16\times 3=0.48.
\]
维持资本密度不变所需投资为
\[
(n+\delta)k_0=(0.01+0.03)\times 27=1.08.
\]

由于
\[
0.48<1.08,
\]
所以人均资本将下降，说明经济位于稳态之上。

由题给更新式，
\[
k_1=\frac{(1-0.03)\times 27+0.16\times 3}{1.01}
=\frac{26.19+0.48}{1.01}
=\frac{26.67}{1.01}\approx 26.41.
\]

稳态满足
\[
s f(k)=(n+\delta)k
\quad\Longrightarrow\quad
0.16k^{1/3}=0.04k.
\]
当 $k>0$ 时，可化为
\[
0.16=0.04k^{2/3}
\quad\Longrightarrow\quad
k^{2/3}=4
\quad\Longrightarrow\quad
k^*=4^{3/2}=8.
\]
\ansmath{k^*=8.}

对于 Cobb--Douglas 形式 $y=k^\alpha$，黄金规则储蓄率为
\[
s_{\text{gold}}=\alpha=\frac13.
\]
\ansmath{s_{\text{gold}}=\frac13.}

因此本题当前储蓄率 $0.16$ 低于黄金规则储蓄率，但由于初始 $k_0=27$ 远高于稳态 $8$，短期内人均资本仍会下降；这说明“是否位于黄金规则”与“是否高于当前稳态”是两个不同判断。
\end{solution}
\end{question}

\begin{question}
有两种风险资产 $A$ 与 $B$，各有两种状态，概率都为 $0.5$。其收益率分别为：

\begin{center}
\begin{tabular}{c|cc}
状态 & 资产 $A$ & 资产 $B$\\ \hline
1 & $30\%$ & $-10\%$\\
2 & $-10\%$ & $30\%$
\end{tabular}
\end{center}

\begin{enumerate}[label=(\arabic*),leftmargin=2.5em]
  \item 分别求资产 $A$、$B$ 的期望收益率和标准差；
  \item 若投资者按 $1:1$ 比例持有 $A$ 与 $B$，求组合的期望收益率和标准差；
  \item 解释这个例子为什么能说明“分散化降低非系统性风险”，以及为什么现实中不意味着风险一定能降为零。
\end{enumerate}

\begin{solution}
对资产 $A$，
\[
E(r_A)=0.5\times 30\%+0.5\times(-10\%)=10\%.
\]
偏差分别为 $20\%$ 与 $-20\%$，故方差
\[
\mathrm{Var}(r_A)=0.5(20\%)^2+0.5(-20\%)^2=0.04,
\]
标准差
\[
\sigma_A=20\%.
\]
同理，
\[
E(r_B)=10\%,\qquad \sigma_B=20\%.
\]

若按 $1:1$ 组合持有，则两种状态下组合收益率都为
\[
r_P=0.5\times 30\%+0.5\times (-10\%)=10\%
\]
或
\[
r_P=0.5\times (-10\%)+0.5\times 30\%=10\%.
\]
因此
\[
E(r_P)=10\%,\qquad \sigma_P=0.
\]
\ansmath{E(r_A)=E(r_B)=10\%,\quad \sigma_A=\sigma_B=20\%,}
\ansmath{E(r_P)=10\%,\quad \sigma_P=0.}

这里两种资产在两个状态下呈完全负相关，个别波动被恰好抵消，所以组合风险降为零。这说明分散化能够消除资产特有风险。但现实中资产通常不会如此完美地负相关，且宏观冲击会同时影响多种资产，因此系统性风险一般无法被完全消除。
\end{solution}
\end{question}

